\documentclass[11pt]{article}
\usepackage[margin=25mm]{geometry}
\usepackage{fontspec}
\setmainfont{Latin Modern Roman}
\usepackage{amsmath,amssymb,mathtools}
\usepackage{unicode-math}
\setmathfont{Latin Modern Math}
\usepackage{xurl,hyperref}
\hypersetup{colorlinks=true,urlcolor=blue}
\setcounter{secnumdepth}{0}
\setlength{\parindent}{0pt}
\setlength{\parskip}{5pt}
\setlength{\emergencystretch}{3em}
\newcommand{\sref}[2]{\hyperlink{source:#1:#2}{[S#2]}}
\newcommand{\cataloguescope}[1]{\paragraph{Recorded target.}#1}
\newcommand{\R}{\mathbb R}
\newcommand{\C}{\mathbb C}
\newcommand{\N}{\mathbb N}
\newcommand{\Z}{\mathbb Z}
\newcommand{\Q}{\mathbb Q}
\newcommand{\F}{\mathbb F}
\newcommand{\eps}{\varepsilon}
\begin{document}
% UnsolvedMath ID 20002555; source catalogue code AIM-PHYSICS-0063
\setcounter{section}{20002554}
\section{Compatibility and reversible shooting for connections between relative equilibria}\label{20002555}
{\small UnsolvedMath ID 20002555 · Mathematical Physics}
\subsection{Definitions and mathematical statement}

\paragraph{Shared notation.}$\mathbb N=\{1,2,\ldots\}$; the standard number fields have their usual meanings. All additional conventions are stated locally.


\paragraph{Definitions and notation.}In the Newtonian $n$-body problem, positions $q_i\in\mathbb R^d$ and masses $m_i>0$ obey $m_i\ddot q_i=\nabla_{q_i}U$, where $U(q)=\sum_{i<j}m_im_j/|q_i-q_j|$ and the gravitational constant is normalized to one. Here $T=\tfrac12\sum_i m_i|\dot q_i|^2$, $E=T-U$, and a collision has $q_i=q_j$. The center of mass can be fixed at zero; moment of inertia is $I=\sum_i m_i|q_i|^2$. A central configuration has $\nabla_{q_i}U=-\lambda m_iq_i$ in center-of-mass coordinates for a common scalar $\lambda$. After fixing $I=1$, it is a critical point of $U$ on the collision-free inertia ellipsoid. Congruent configurations and scales are factored out for a finiteness count. A relative equilibrium is a rigidly rotating solution with time-independent shape. A homoclinic orbit tends to the same invariant equilibrium orbit in forward and backward time; a heteroclinic orbit tends to two different ones. For relative equilibria these limits are naturally described after rotation reduction or in a compatible rotating frame. Their stable and unstable manifolds must be compared on common conserved-energy and angular-momentum levels.

\paragraph{Catalogue wording retained for traceability.}
Problem 26. Existence proof for homoclinic and heteroclinic orbits between unstable rela-tive equilibrium solutions. (Rick Moeckel)
\paragraph{Scope and source qualification.}The source does not quantify over every unstable relative equilibrium, fix masses or prescribe a particular pair. Existence in a specified compatible setting is required, and a universal connection claim would be stronger than the source. The catalogue title/status was paired with a machine-generated research attempt; that attempt is not peer-reviewed evidence of settlement. The statement covers the original source questions rather than only the title’s subcase.
\subsection{Short English statement}
Prove homoclinic and heteroclinic connections involving unstable relative equilibria in celestial mechanics.
\subsection{Research attempt}
\paragraph{Model, date and investigated scope.}
OpenAI GPT-6 Astra (configured model \texttt{gpt-6-astra}), ultra reasoning; generated 11 October 2026. Derived scale-invariant compatibility of relative equilibria and proved a reversible homoclinic shooting implication. No actual n-body homoclinic or heteroclinic intersection is constructed.

\paragraph{Sources and literature check.}
Read the supplied complete definition, catalogue wording, scope qualifications and source list for this canonical ID. The original-source citations are inherited as supplied references; no fresh primary-source or novelty search was performed for this attempt. No current open-status claim is inferred from this calculation.

\paragraph{A conserved-quantity compatibility test.}For a planar circular relative equilibrium generated by a central shape, let its inertia and potential at a chosen scale be \(I,U\). Force balance gives \(\omega^2=U/I\). Therefore
\[C=\omega I,\quad E=-U/2,\quad C^2E=-IU^2/2.
\]
The last quantity is independent of scale, since \(I\) scales quadratically and \(U\) inversely. If a heteroclinic trajectory tends in phase space, after an appropriate reduction, to two relative equilibria, continuity and conservation require their energy and angular momentum to agree. In particular their shape values \(IU^2\) must match. Instability by itself does not supply this compatibility.
\paragraph{A reversible shooting criterion.}Let a smooth flow \(\Phi_t\) admit an involution \(R\) with \(R\Phi_t=\Phi_{-t}R\), and let \(p\) be an R-fixed equilibrium in the reduced space. If a point \(x\ne p\) lies in \(W^u(p)\cap\operatorname{Fix}R\), then it also lies in \(W^s(p)\): as \(t\to+\infty\), \(\Phi_t x=R\Phi_{-t}x\to Rp=p\). Hence its complete trajectory is homoclinic, assuming it exists without collision for all time. This proves the implication from a shooting intersection, not existence of that intersection.
\paragraph{Gap.}No concrete n-body masses, common level, unstable branch or validated intersection is exhibited. The first missing step is establishing such an intersection while excluding collision or noncompact escape. Heteroclinic connections additionally require the compatible target pair, and the source does not ask for all unstable equilibria indiscriminately. The calculation is a necessary-condition and conditional-shooting benchmark, not a connection-existence theorem.

\paragraph{Verification and adversarial audit.}
Checked circular energy/angular-momentum factors, scale invariance, continuity of integrals and reversible-time limit argument with completeness hypothesis. This was an LLM adversarial audit, not independent human verification; no theorem prover was used. Prompts A then B applied to this ID: Checked circular energy/angular-momentum factors, scale invariance, continuity of integrals and reversible-time limit argument with completeness hypothesis. Scope remains limited: Find and rigorously validate stable/unstable intersections on common conserved levels while controlling collisions and escape.

Coordinator checked scale invariance of the relative-equilibrium energy/angular-momentum condition and the reversible homoclinic shooting implication. This additional check was model-assisted, not independent human verification.

\paragraph{Tools, process and contributions.}
Python 3 and local shell used to read supplied files and serialize this manuscript; mathematics derived symbolically in the text, with no numerical experiment or formal proof assistant. Contribution prompt A was followed by prompt B. GPT-6 Astra with ultra reasoning drafted the research and performed its stated self-audit; any additional coordinator contribution is identified in the verification and audit. A GPT-6-based Codex coordinator assembled the manuscript and checked repository integration; its exact variant and reasoning configuration were not exposed. No human mathematical contribution, independent human proof check or formal proof certification is supplied. Complete derivations and the consulted source references are included here; private reasoning transcripts are not published.

\paragraph{Final status and first remaining gap.}
\textbf{UNRESOLVED}. Find and rigorously validate stable/unstable intersections on common conserved levels while controlling collisions and escape.

\section{Completion Estimate}
\noindent\textbf{COMPLETION: 0\%}
The calculation is an elementary benchmark or conditional reduction; it establishes no claimed new portion of the complete source target.

\subsection{Sources}
\begin{itemize}
\item[\hypertarget{source:20002555:1}{S1}] ULAM.AI and the UnsolvedMath Contributors, UnsolvedMath, record 20002555 (AIM-PHYSICS-0063). Catalogue curation under CC BY 4.0; definitions and English exposition added locally.
\url{https://huggingface.co/datasets/ulamai/UnsolvedMath}.
\item[\hypertarget{source:20002555:2}{S2}] American Institute of Mathematics, Variational Methods in Celestial Mechanics; notes by Kuo-Chang Chen, version 22 June 2003, source item 26.
\url{https://aimath.org/WWN/varcelest/varcelest.pdf}.
\end{itemize}
\label{end:20002555}
\end{document}
